% ============================================================
% madXfinder — 시행착오 기록: 우리가 어떻게 알아냈나
% Overleaf: New Project → Upload → 이 파일.
% ★ Compiler: XeLaTeX 필수 (Menu → Compiler). pdfLaTeX은 한글 폰트 처리로 시간초과 남
% 상세 근거·수치: docs/KJH/의사결정-기록.md
% ============================================================
\documentclass[11pt]{article}
\usepackage{kotex}
\usepackage{amsmath}
\usepackage{booktabs}
\usepackage{geometry}
\geometry{a4paper, margin=25mm}
\usepackage{hyperref}
\hypersetup{colorlinks=true, linkcolor=blue, urlcolor=blue}
\usepackage{fancyvrb}

\title{madXfinder 시행착오 기록:\\로블록스 API의 한계를 실측으로 규명하기까지}
\author{madcamp 26s-w1-c3-09}
\date{}

\begin{document}
\maketitle

\begin{abstract}
로블록스는 API 제한 정책을 공개하지 않는다. 우리는 추천 서비스를 만들기 위해
그 한계를 직접 측정해야 했고, 그 과정에서 세 번 결론이 뒤집혔다.
이 문서는 ``무엇을 만들었나''가 아니라 \textbf{``어떻게 알아냈나''}의 기록이다.
잘못 짚었던 것, 측정 도구가 만든 착시, 그리고 그것을 어떻게 바로잡았는지를 남긴다.
\end{abstract}

\section{요약: 세 번의 반전}

\begin{center}
\begin{tabular}{clll}
\toprule
 & 그때의 결론 & 왜 그렇게 믿었나 & 뒤집은 것 \\
\midrule
1차 & ``동시성 N까지 허용'' & 동시 요청 수를 늘리며 테스트 & 같은 N인데 결과가 널뜀 \\
2차 & ``IP 전역 공유 한도'' & 두 API를 동시에 치니 서로 무너짐 & 측정 도구의 착시였음 \\
3차 & \textbf{``엔드포인트별 독립''} & 전 API 재측정 + 조합 실험 & (최종 --- 예외 1개 발견) \\
\bottomrule
\end{tabular}
\end{center}

\section{1차 시도: ``동시에 몇 개까지 되나?'' --- 잘못된 질문}

처음엔 다들 하듯 동시 요청 수를 늘려봤다. 그런데 이상했다:
\textbf{동시성 20이 어떤 때는 에러 0\%, 어떤 때는 94\%.} 같은 조건인데 결과가 다르다면
변수를 잘못 잡은 것이다.

원인은 \textbf{토큰버킷}이었다. 로블록스는 ``동시에 몇 개''가 아니라
``통(bucket)에 토큰이 몇 개 남았나''로 제한한다. 통이 가득할 땐 동시성 20도 통과하고,
직전 테스트가 통을 비워놨으면 같은 20이 폭발한다. 순간 동시성은 애초에 측정 대상이
아니었다 --- \textbf{질문을 ``초당 몇 개(rate)''로 바꿔야 했다.}

\begin{Verbatim}[frame=single, fontsize=\small]
교훈 1. 결과가 재현되지 않으면 측정값이 아니라 "측정 모델"을 의심하라.
\end{Verbatim}

\section{2차 시도: ``IP 전역 한도''라는 착시}

rate 기준으로 다시 재니 값이 나오기 시작했다. 그런데 서로 다른 두 도메인
(그룹 API + 게임 API)을 동시에 부하하자 둘 다 무너졌다. ``아, 도메인이고 뭐고
IP당 전체 한도(약 13회/초)를 공유하는구나''라고 결론 냈다.

\textbf{이것이 착시였다.} 당시 테스트는 스레드 100개 + 동기 HTTP로 짰는데,
우리 쪽 PC의 소켓이 고갈되면서 연결 에러가 났던 것이다. 로블록스가 막은 게 아니라
\textbf{우리 측정 도구가 스스로 무너진 것}을 상대방 탓으로 읽었다.

비동기(asyncio) 기반으로 도구를 갈아엎고 다시 재자 그림이 완전히 달라졌다.

\begin{Verbatim}[frame=single, fontsize=\small]
교훈 2. 부하 테스트에서 에러가 나면, 상대가 막은 건지
        내 도구가 못 버틴 건지부터 구분하라. (연결 에러 != 429)
\end{Verbatim}

\section{3차 시도: 전수 측정 --- ``엔드포인트별 독립 버킷''}

이번엔 가정을 버리고 전부 쟀다. 5개 도메인 13개 엔드포인트를 각각 단독으로,
그리고 조합으로.

\subsection{단독 측정 (지속 가능 rate, 회/초)}
\begin{center}
\small
\begin{tabular}{llll}
\toprule
빠름 & \multicolumn{3}{l}{그룹멤버 7.3 · 아이콘 7.3 · 미디어 7.2 · 투표 6.9 · 추천 6.2 · 즐겨찾기 6.1} \\
중간 & \multicolumn{3}{l}{검색 2.8 · 장소 1.6 · 썸네일 1.4 · 차트 1.4} \\
\textbf{병목} & \multicolumn{3}{l}{\textbf{닉네임 1.2 · 게임상세 1.1 · 그룹정보 0.3}} \\
\bottomrule
\end{tabular}
\end{center}

엔드포인트마다 값이 제각각이다 --- 이미 ``전역 한도''와는 안 맞는 그림.

\subsection{결정적 실험: 조합해도 각자 유지되나}

같은 도메인(games) 안의 세 엔드포인트를 \textbf{동시에} 각각 초당 2회로 쳤다:
\begin{Verbatim}[frame=single, fontsize=\small]
즐겨찾기 2/s + 추천 2/s + 상세 2/s  →  합계 6/s, 전원 무사

만약 도메인 공유라면: 합계가 한도에 눌려 각자 떨어졌어야 함
실제: 각자 그대로 → 엔드포인트마다 독립된 통이 있다!
→ 서로 다른 API를 조합하면 총 처리량이 "합산"된다
\end{Verbatim}

이 발견이 아키텍처를 결정했다: 호출 예산 관리는 전역 하나가 아니라
\textbf{엔드포인트별로} 해야 하고, 여러 API를 섞어 쓰면 처리량이 배가된다.

\subsection{예외 하나: 유저 조회 계열은 통을 공유한다}

닉네임 변환에는 경로가 3개 있다 (단건 GET, 대량 POST 2종). 이들을 조합 실험하니
\textbf{이번엔 합산되지 않고 나눠 가졌다} --- 유저 조회 계열만은 하나의 통을 공유한다.
``전부 독립''이라고 일반화할 뻔한 것을 조합 실험이 잡아냈다.

\begin{Verbatim}[frame=single, fontsize=\small]
교훈 3. "패턴을 찾았다"고 생각한 순간에도 예외 검증을 하라.
        규칙(독립)과 예외(users 공유)는 실험 하나 차이였다.
\end{Verbatim}

\subsection{한 통에 몇 개까지 담아 보낼 수 있나 (배치 상한)}

일부 API는 한 호출에 ID 여러 개를 받는다. 개수를 늘려가며 거부(400)가 시작되는
지점을 이분탐색으로 확정했다:

\begin{center}
\begin{tabular}{lll}
\toprule
API & 호출당 최대 & 실효 처리량 \\
\midrule
게임 상세 & 50개 (51부터 거부) & $1.1 \times 50 \approx$ 초당 55게임 \\
투표·아이콘·썸네일 & 각 100개 & 최대 초당 수백 게임 \\
\bottomrule
\end{tabular}
\end{center}

``초당 1.1회''라던 최악의 병목(게임 상세)이, 묶어 보내니 \textbf{초당 55게임}이 됐다.
rate limit은 ``호출당''이지 ``게임당''이 아니었던 것이다.

\section{429의 대가: 밀어붙이면 진다}

한도를 넘겨 429가 뜬 뒤에도 계속 찌르면, 회복이 30초 이상 걸리는 페널티가 쌓였다.
반대로 429를 만나면 즉시 속도를 절반으로 줄이고 물러나는 방식(AIMD)은
장시간 안정적으로 한도 근처를 유지했다. \textbf{``빠르게 많이''보다 ``느려도 꾸준히''가
총량에서 이긴다} --- 그래서 모든 수집기는 AIMD + 여유 마진으로 돈다.

\section{수집 전략의 시행착오}

\subsection{누구의 즐겨찾기를 모을 것인가}
게임의 팬을 찾는 경로로 ``그 게임 그룹의 멤버''를 쓰는데, 어떤 순서로 도느냐가
품질을 갈랐다:

\begin{center}
\begin{tabular}{lll}
\toprule
순회 방향 & 즐겨찾기 보유 & 평균 개수 \\
\midrule
최신 가입순 (Desc) & 낮음 & 6.6개 (막 유입된 뉴비) \\
\textbf{오래된 가입순 (Asc)} & \textbf{95--98\%} & \textbf{27--30개} (초기부터 판 코어층) \\
\bottomrule
\end{tabular}
\end{center}

덤으로 Asc는 과거 가입자 순서가 고정이라 \textbf{커서가 밀리지 않는다} ---
200명에서 끊고 저장해두면 나중에 201번째부터 정확히 이어받는다.

\subsection{버린 똑똑한 아이디어들}
\begin{itemize}
  \item \textbf{뉴비/개발자 스킵}: 앞 100명을 건너뛰고 수집한 결과와 비교 실험 →
        추천 상위가 완전히 동일. 개발자 계정은 앞쪽에 몰려있지도 않았다. 복잡도만 늘어 폐기.
  \item \textbf{코어팬 임계값}: ``즐겨찾기 N개 이상인 진짜 팬만 저장''하려 했으나,
        호출 예산이 귀해서 어렵게 조회한 유저를 버리는 것 자체가 낭비 → 무조건 저장으로.
  \item \textbf{게임당 표본 키우기}: 100명 단위로 늘려보니 ``새 게임 발견''은 500명쯤에서
        포화. 200명이면 상위 취향은 이미 안정적 → 한 게임을 깊게보다 \textbf{많은 게임을
        200명씩 넓게}가 예산 대비 효율적.
\end{itemize}

\section{알고리즘 검증: 장르가 다른 두 게임으로}

양산형 시뮬레이터와 애니메이션 IP 게임 --- 성격이 정반대인 두 게임의 팬 200명씩으로
co-favorite 추천을 뽑아 교차 검증했다. 각자의 추천 상위에 각자 성격의 게임들이
정확히 정렬됐고, 이때 배운 것들:

\begin{itemize}
  \item \textbf{장르 필드는 필요 없다.} 겹침 신호가 장르를 ``스스로'' 반영한다.
  \item \textbf{200명 중 7명(3.5\%)도 신호다.} 처음엔 노이즈로 버리려 했으나,
        저겹침 게임이 오히려 숨은 취향작인 경우가 많았다 --- 발견 섹션의 재료.
  \item \textbf{유명도 보정은 세기 조절이 전부다.} 약하면 대작만 나오고, 세면
        이상한 소형작이 튄다. 하나의 값으로 안 되니 두 섹션(인기/발견)으로 분리.
  \item \textbf{죽은 게임은 동접으로 거른다.} 방문수 필터로는 폐쇄된 게임이 남았다.
\end{itemize}

\section{시간의 한계: ``실시간 서비스''를 포기한 이유}

처음엔 유저 요청마다 로블록스를 실시간으로 부를 생각이었다. 측정 결과가 그 계획을
부쉈다:

\begin{Verbatim}[frame=single, fontsize=\small]
게임 상세 = 초당 1.1회 (단건 기준)
→ 유저 10명이 동시에 게임을 클릭하면 마지막 사람은 ~9초 대기
→ 100명이면 90초. 서비스가 아니다.
\end{Verbatim}

결론: 실시간으로 부를 수 있는 건 \textbf{딱 두 가지}로 압축됐다 ---
닉네임$\to$ID 변환(신규 유저 1회)과 그 유저의 즐겨찾기 조회(1회).
나머지 전부(게임 정보, 썸네일, 추천 계산)는 배치가 미리 DB에 만들어 둔다.
인기 상위 게임부터 차례로 캐싱하면(top-down), 추천에 등장하는 게임은 인기 편향이라
대부분 캐시에 이미 있다.

\section{교훈 모음}

\begin{enumerate}
  \item 재현 안 되는 측정은 측정이 아니다 --- 모델(동시성 vs rate)부터 다시.
  \item 부하 테스트의 에러는 절반이 내 도구 탓이다 (스레드 소켓 고갈의 착시).
  \item 규칙을 찾았어도 예외를 찾는 실험을 한 번 더 (users 공유 버킷).
  \item 한도를 밀어붙이면 페널티로 되갚음 당한다 --- 물러나는 쪽이 총량에서 이긴다.
  \item ``호출당'' 제한은 묶음(배치)으로 우회된다 (1.1회/초 = 55게임/초).
  \item 어렵게 얻은 데이터는 버리지 않는 구조로 설계한다 (무조건 저장).
  \item 실측이 아키텍처를 결정하게 하라 --- 우리의 ``배치 중심 + 실시간 2호출'' 구조는
        취향이 아니라 측정값이 강제한 것이다.
\end{enumerate}

\bigskip
\noindent 모든 수치의 원본 기록과 결정 근거: \texttt{docs/KJH/의사결정-기록.md}

\end{document}
